% GENERATED FILE. Edit canonical lessons, then run npm run generate:books.
% canonical-source-hash: fnv1a64:603a9d0933876eca
% canonical-lessons: PT-C327-rei, PT-C327-rainha, PT-C327-reino, PT-C327-trono, PT-C327-coroa

\chapter{Reis e rainhas}
\label{ch:pt-reis-e-rainhas}
\hlchaptermodality{327}

\begin{chapteropening}
\textbf{By the end of this chapter:} \emph{I can talk about kings and queens, kingdoms, thrones and crowns.}

Complete the last lesson of chapter 327: I can talk about kings and queens, kingdoms, thrones and crowns.
\end{chapteropening}

\section[o rei]{o rei — the king}

\label{lesson:PT-C327-rei}



\begin{quote}
Before the new one: say the Portuguese for the fable, then the Portuguese for the fairy.
\end{quote}

\subsection*{You'll want to know: o rei}
\textbf{o rei} — \textquotedblleft{}the king\textquotedblright{}.

\begin{quote}
O rei vivia num castelo grande.
\end{quote}

\begin{quote}
\emph{The king lived in a big castle.}
\end{quote}

\emph{Viver como um rei} — to live like a king.

\begin{cousinweb}
\emph{rei} is inherited from Latin \emph{rēx, rēgem} \textquotedblleft{}a king\textquotedblright{} (EtymDB).
\end{cousinweb}

\begin{grammarlens}[title={What you've built}]
One more word for telling, explaining and describing at B1.
\end{grammarlens}

\subsection*{Your turn}
\begin{itemize}
\raggedright
  \item \emph{Say it:} \emph{o rei}
  \item \emph{Say it:} \emph{o rei}, once more
  \item \emph{Say it:} a sentence of your own with \emph{o rei}
  \item \emph{Recall:} say the Portuguese for the fable, then the Portuguese for the fairy, then say \emph{o rei} again
\end{itemize}

\begin{tcolorbox}[breakable,colback=teal!4,colframe=teal!35!black,title={Before you move on}]
What does o rei mean? (\textquotedblleft{}The king\textquotedblright{}.) Say it once more.
\end{tcolorbox}
\section[a rainha]{a rainha — the queen}

\label{lesson:PT-C327-rainha}



\begin{quote}
Before the new one: say the Portuguese for the fairy, then the Portuguese for the king.
\end{quote}

\subsection*{You'll want to know: a rainha}
\textbf{a rainha} — \textquotedblleft{}the queen\textquotedblright{}.

\begin{quote}
A rainha usava um vestido azul.
\end{quote}

\begin{quote}
\emph{The queen wore a blue dress.}
\end{quote}

Also the queen in chess.

\begin{cousinweb}
\emph{rainha} comes down from Old Portuguese \emph{raỹa} (EtymDB), from Latin \emph{rēgīna} \textquotedblleft{}a queen\textquotedblright{}, from \emph{rēx} \textquotedblleft{}a king\textquotedblright{}.
\end{cousinweb}

\begin{grammarlens}[title={What you've built}]
One more word for telling, explaining and describing at B1.
\end{grammarlens}

\subsection*{Your turn}
\begin{itemize}
\raggedright
  \item \emph{Say it:} \emph{a rainha}
  \item \emph{Say it:} \emph{a rainha}, once more
  \item \emph{Say it:} a sentence of your own with \emph{a rainha}
  \item \emph{Recall:} say the Portuguese for the fairy, then the Portuguese for the king, then say \emph{a rainha} again
\end{itemize}

\begin{tcolorbox}[breakable,colback=teal!4,colframe=teal!35!black,title={Before you move on}]
What does a rainha mean? (\textquotedblleft{}The queen\textquotedblright{}.) Say it once more.
\end{tcolorbox}
\section[o reino]{o reino — the kingdom}

\label{lesson:PT-C327-reino}



\begin{quote}
Before the new one: say the Portuguese for the king, then the Portuguese for the queen.
\end{quote}

\subsection*{You'll want to know: o reino}
\textbf{o reino} — \textquotedblleft{}the kingdom\textquotedblright{}.

\begin{quote}
O reino era pequeno, mas rico.
\end{quote}

\begin{quote}
\emph{The kingdom was small but rich.}
\end{quote}

\emph{O reino animal} — the animal kingdom.

\begin{cousinweb}
\emph{reino} goes back to Latin \emph{rēgnum} \textquotedblleft{}royal power, a kingdom\textquotedblright{} (EtymDB), from \emph{rēx} \textquotedblleft{}a king\textquotedblright{}.
\end{cousinweb}

\begin{grammarlens}[title={What you've built}]
One more word for telling, explaining and describing at B1.
\end{grammarlens}

\subsection*{Your turn}
\begin{itemize}
\raggedright
  \item \emph{Say it:} \emph{o reino}
  \item \emph{Say it:} \emph{o reino}, once more
  \item \emph{Say it:} a sentence of your own with \emph{o reino}
  \item \emph{Recall:} say the Portuguese for the king, then the Portuguese for the queen, then say \emph{o reino} again
\end{itemize}

\begin{tcolorbox}[breakable,colback=teal!4,colframe=teal!35!black,title={Before you move on}]
What does o reino mean? (\textquotedblleft{}The kingdom\textquotedblright{}.) Say it once more.
\end{tcolorbox}
\section[o trono]{o trono — the throne}

\label{lesson:PT-C327-trono}



\begin{quote}
Before the new one: say the Portuguese for the queen, then the Portuguese for the kingdom.
\end{quote}

\subsection*{You'll want to know: o trono}
\textbf{o trono} — \textquotedblleft{}the throne\textquotedblright{}.

\begin{quote}
O rei sentou-se no trono.
\end{quote}

\begin{quote}
\emph{The king sat down on the throne.}
\end{quote}

\emph{Subir ao trono} — to come to the throne.

\begin{cousinweb}
\emph{trono} comes down from Old Portuguese (EtymDB) and goes back to Latin \emph{thronus}, from Greek \emph{thrónos} \textquotedblleft{}a seat, a chair\textquotedblright{} (Lewis \& Short).
\end{cousinweb}

\begin{grammarlens}[title={What you've built}]
One more word for telling, explaining and describing at B1.
\end{grammarlens}

\subsection*{Your turn}
\begin{itemize}
\raggedright
  \item \emph{Say it:} \emph{o trono}
  \item \emph{Say it:} \emph{o trono}, once more
  \item \emph{Say it:} a sentence of your own with \emph{o trono}
  \item \emph{Recall:} say the Portuguese for the queen, then the Portuguese for the kingdom, then say \emph{o trono} again
\end{itemize}

\begin{tcolorbox}[breakable,colback=teal!4,colframe=teal!35!black,title={Before you move on}]
What does o trono mean? (\textquotedblleft{}The throne\textquotedblright{}.) Say it once more.
\end{tcolorbox}
\section[a coroa]{a coroa — the crown}

\label{lesson:PT-C327-coroa}



\begin{quote}
Before the new one: say the Portuguese for the kingdom, then the Portuguese for the throne.
\end{quote}

\subsection*{You'll want to know: a coroa}
\textbf{a coroa} — \textquotedblleft{}the crown\textquotedblright{}.

\begin{quote}
A coroa da rainha era de ouro.
\end{quote}

\begin{quote}
\emph{The queen’s crown was made of gold.}
\end{quote}

Also the top of the head, and the heads side of a coin.

\begin{cousinweb}
\emph{coroa} comes down from Old Portuguese \emph{corõa} (EtymDB), from Latin \emph{corōna} \textquotedblleft{}a wreath, a crown\textquotedblright{}, from Greek \emph{kor\'{\={o}}nē} \textquotedblleft{}something curved\textquotedblright{}.
\end{cousinweb}

\begin{grammarlens}[title={What you've built}]
One more word for telling, explaining and describing at B1.
\end{grammarlens}

\subsection*{Your turn}
\begin{itemize}
\raggedright
  \item \emph{Say it:} \emph{a coroa}
  \item \emph{Say it:} \emph{a coroa}, once more
  \item \emph{Say it:} a sentence of your own with \emph{a coroa}
  \item \emph{Recall:} say the Portuguese for the kingdom, then the Portuguese for the throne, then say \emph{a coroa} again
\end{itemize}

\begin{tcolorbox}[breakable,colback=teal!4,colframe=teal!35!black,title={Before you move on}]
What does a coroa mean? (\textquotedblleft{}The crown\textquotedblright{}.) Say it once more.
\end{tcolorbox}
